\documentclass[14pt]{extreport}
\usepackage{graphicx}
\usepackage{gost}
\usepackage[T1,T2A]{fontenc}
\usepackage[utf8]{inputenc}
\PassOptionsToPackage{russian,english}{babel}
\usepackage[english,russian]{babel}
\usepackage{tempora}
\usepackage{hyperref}
\linespread{1.3}
\setlength{\parindent}{1.25cm}
\usepackage{hyperref}

\usepackage{multirow}
\usepackage{makecell}
\usepackage{longtable}
\usepackage{array}
\newcolumntype{P}[1]{>{\raggedright\arraybackslash}p{#1}}

\usepackage{fancyhdr}
\pagestyle{fancy}
\fancyhf{} 
\fancyhead{} 
\fancyfoot[C]{\thepage} 

\usepackage{amsmath, amssymb}
\usepackage{tcolorbox}
\usepackage{enumitem}
\usepackage{geometry}
\usepackage{amsthm}
\tcbuselibrary{breakable}

\newtcolorbox{thbox}[1][]{colback=orange!10, colframe=orange!60, title=Теорема, #1, breakable=true, toprule at break=0pt, % верхняя рамка не показывается при разрыве
    bottomrule at break=0pt}
\newtcolorbox{prbox}[1][]{colback=green!5, colframe=green!60!gray, title=Доказательство, breakable=true, #1, toprule at break=0pt, % верхняя рамка не показывается при разрыве
    bottomrule at break=0pt}
\newtcolorbox{defbox}[1][]{colback=blue!5, colframe=blue!60, title=Определение, #1, breakable=true, toprule at break=0pt, % верхняя рамка не показывается при разрыве
    bottomrule at break=0pt}
\newtcolorbox{impbox}[1][]{colback=red!5, colframe=red!60, title=Важно, #1, breakable=true, toprule at break=0pt, % верхняя рамка не показывается при разрыве
    bottomrule at break=0pt}

\newcommand{\defeq}[0]{\stackrel{\text{def}}{=}}

\begin{document}
\thispagestyle{empty}
	\title{Высшая алгебра. Экзамен.}
	\maketitle

	\newpage
	\tableofcontents

    \newpage
    \chapter{Линейные операторы. Общие понятия.}

    \begin{defbox}
        Пусть есть пространства $U$ и $V$. \textbf{Линейным оператором }A $U \stackrel{\text{A}}{\rightarrow} V$ называется линейное отображение $U$ в $V$, такое, что:
        \begin{enumerate}
            \item $\forall x \in U \: \forall y \in V \quad A(x+y)=A(x)+A(y)$
            \item $\forall \alpha \in \mathbb{C}(\mathbb{R}) \quad A(\alpha x)=\alpha A(x)$
        \end{enumerate}
    \end{defbox}
    \begin{thbox}
        Действие любого оператора задается действием на базисные вектора
    \end{thbox}
    \begin{prbox}
        $$
        X = \begin{pmatrix}
            x_1\\
            x_2\\
            \vdots\\
            x_n
        \end{pmatrix} = x_1e_1 + \cdots + x_ne_n
        $$
        $$
        AX = A(x_1e_1 + \cdots + x_ne_n) = A(x_1e_1) + \cdots + A(x_ne_n) = 
        $$
        $$
        =x_1Ae_1 + \cdots + x_nAe_n
        $$

        Получается: 
        $$
        e_1 \cdots e_n \:\text{независимы} \Leftrightarrow Ae_1\cdots Ae_n \: \text{независимы}
        $$

        
    \end{prbox}
    \begin{impbox}
        Матрица оператора в заданном базисе представляет собой столбцы-векторы, являющиеся результатом действия оператора на базисные векторы
    \end{impbox}

    \chapter{Собственные числа и собственные векторы.}
    \begin{defbox}
        Число $\lambda$ называется \textbf{собственным числом} оператора $A$, если существует ненулевой вектор такой, что $Ax = \lambda x$
    \end{defbox}
    \begin{thbox}
        Число $\lambda$ называется \textbf{собственным числом} оператора $A$ тогда и только тогда, когда $\det(A - \lambda E) = 0$
    \end{thbox}
    \begin{prbox}
        1. $x$ -- собственный $\Leftrightarrow \det (A - \lambda E) = 0$ 
        $$
        Ax = \lambda x = \lambda E x
        $$
        $$
        (A - \lambda E )x = 0
        $$
        Где $A - \lambda E$ -- линейный оператор
    \end{prbox}
\end{document}